\origpage{57}{GRAPH MINORS. IX}
\hspace*{12pt}(7) $\Omega(a, b) = \emptyset$.

For suppose that $c \in \Omega(a, b)$. Let $G'$ be the component of $G$ containing $c$. If $G' \cap P$ is not null, let $P_1$ be a $\{c\} \to V(G' \cap P)$ path in $G'$; then $P_1 \subseteq G_1$, since no internal vertex of $P_1$ is in $G_1 \cap G_2$ and $c \in V(G_1) - V(G_2)$. Hence there is an $\Omega$-path $P'$ with $P' \subseteq G_1$ and one end $c$, contrary to (6). Thus $G' \cap P$ is null, and so $G' \subseteq G_1$. By (6), there is no $\Omega$-path $P'$ with $P' \subseteq G'$, and so only one vertex of $G'$ is in $\bar{\Omega}$. But $(G, \Omega)$ is 3-connected, and so $V(G') = \{c\}$, contrary to (2). This verifies (7).

Let $H = G_2 \setminus V(G_1)$. By the minimality of $G_2$, $\Omega(b, a)$ catches $H$, and $G_1$ is an induced subgraph of $G$.

\hspace*{12pt}(8) \textit{If $H'$ is a subgraph of $G$ such that $H \subseteq H'$, $\Omega(b, a)$ catches $H'$, and there is a path of $G$ between $a$ and $b$ avoiding $V(H')$, then $H' = H$.}

For let $P'$ be a path of $G$ between $a$ and $b$ avoiding $V(H')$. Then $P' \subseteq G_1$, since $V(P') \subseteq V(G_1)$ and $G_1$ is an induced subgraph of $G$. Let $\mathcal{P}' = (\mathcal{P} - \{P\}) \cup \{P'\}$; then $\mathcal{P}'$ is a transaction satisfying (a) and (b) (since $\Omega(b, a) \subseteq V(H)$ and $\Omega(a, b) = \emptyset$). By (5), $C(\mathcal{P}) \subseteq C(\mathcal{P}')$, and so by (c), $C(\mathcal{P}) = C(\mathcal{P}')$. Suppose that $H \ne H'$. Now $H$ is an induced subgraph of $G$, and so there is a vertex $v$ of $H'$ which is not a vertex of $H$. Since $v \in V(H')$ and $\Omega(b, a)$ catches $H'$, there is a $\{v\} \to \Omega(b, a)$ path $Q$ which avoids $V(P')$. Since $v \notin V(H)$ we have $v \in V(G_1)$, and so $P \cap Q$ is non-null; and in particular it follows that we may choose $v$ and $Q$ so that $v \in V(P)$ and $Q$ is a $V(P) \to \Omega(b, a)$ path. Now $v \in V(P)$ and so $v \notin C(\mathcal{P}) = C(\mathcal{P}')$. We deduce that $V(H') \not\subseteq C(\mathcal{P}')$ and so $\mathcal{P}'$ has exactly one other peripheral member $P'' \ne P'$. Let $Q'$ be the minimal subpath of $Q$ between $\{v\}$ and $V(P'') \cup \Omega(b, a)$, and let $v, u$ be the ends of $Q'$. Then $Q'$ meets no non-peripheral member of $\mathcal{P}'$, because $v \notin C(\mathcal{P}')$. Let $P''$ have ends $a'', b''$ where $(a, b, b'', a'')$ is clockwise. Then $u \notin \Omega(a'', b'')$, because $v \notin C(\mathcal{P}')$. Thus either $u \in \Omega(b'', a'')$ or $u \in V(P'')$. Hence any $p - 1$ members of $\mathcal{P} - \{P, P''\}$, together with a suitable path $R$ with $Q' \subseteq R \subseteq P \cup Q' \cup P''$, constitute a crooked transaction with cardinality $p$ in $(G, \Omega)$, which is a contradiction. We deduce that $H' = H$, as claimed. This proves (8).

As $H = G_2 \setminus V(G_1)$ and $G_1$ is induced, it follows that $G_1 = G \setminus V(H)$. Hence, if $G_1 \ne P$, then by (2.1) there would be a separation $(J_1, J_2)$ of $G$ with $\{a, b\} \subseteq V(J_1)$, $|V(J_1 \cap J_2)| = 2$, $H \subseteq J_1$ and $V(J_1) \ne V(G)$. But then $(J_2, J_1)$ would by (7) be a separation of $(G, \Omega)$, contradicting the 3-connectedness of $(G, \Omega)$. It follows that $G_1 = P$.

These results can be summarized as follows.

\hspace*{12pt}(9) \textit{If $P \in \mathcal{P}$ is peripheral and separating, with ends $a, b$ where $\Omega(a, b) \cap V(\mathcal{P}) = \emptyset$, then $\Omega(a, b) = \emptyset$, $P$ is induced, and $\Omega(b, a)$ catches $G \setminus V(P)$.}
